\section{\filetotitle}\label{sec:quantum_amplitude_propagation_in_mixed_themes}
Under the \acrfull{boa}, the total molecular wavefunction, $\ket{\Psi(t)}$, is separable into electronic and nuclear parts because the nuclear masses greatly exceed those of the electrons. However, when two or more electronic states approach degeneracy, the \acrshort{boa} separation fails and nonadiabatic effects become significant. \acrfull{tsh} is a general framework that treats nuclei and electrons separately: nuclei are propagated classically on an adiabatic \acrfull{pes}, while electrons are treated quantum mechanically.

In \acrshort{tsh}, the nuclei follow Newton's equations of motion on a single adiabatic \acrshort{pes} at any instant:
\begin{equation}
\symbf{M}\ddot{\symbf{R}}(t)=-\nabla_{\symbf{R}}E_I,
\end{equation}
where $\symbf{M}$ is the diagonal matrix with masses for each coordinate, $\symbf{R}(t)$ is nuclear geometry at time $t$, and $E_I$ is the electronic energy of the state $I$ at the nuclear configuration $\symbf{R}(t)$. Meanwhile, the electrons evolve according to the \acrshort{tdse}, with the electronic wavefunction $\ket{\Phi(t)}$ expanded in the instantaneous adiabatic eigenvectors of the electronic Hamiltonian $H_{\symrm{el}}$ as
\begin{equation}
\ket{\Phi(t)}=\sum_I c_I(t)\ket{\phi_I},
\end{equation}
where $\phi_I$ satisfies the \acrfull{tise}
\begin{equation}
H_{\symrm{el}}\ket{\phi_I}=E_I\ket{\phi_I},
\end{equation}
and the complex coefficients $c_I$ encode the instantaneous probability amplitudes for occupying each adiabatic state. For notational brevity, we will often suppress explicit dependence on the variables where no ambiguity arises.

Substituting this ansatz into the full electronic \acrshort{tdse},
\begin{equation}
\i\frac{\partial}{\partial t}\ket{\Phi(t)}=H_{\symrm{el}}\ket{\Phi(t)},
\end{equation}
and projecting onto a particular adiabatic state $\ket{\phi_I}$, we obtain a set of coupled equations for the coefficients $c_I$:
\begin{equation}\label{eq:tsh_eom}
\i\frac{\symrm{d}c_I}{\dt}=\sum_J\left(H_{IJ}^{\symrm{el}}-\i\Braket{\phi_I\vert\frac{\partial\phi_J}{\partial t}}\right)c_J,
\end{equation}
where $\sigma_{IJ}=\Braket{\phi_I\vert\frac{\partial\phi_J}{\partial t}}$ is the \acrfull{tdc}. In the purely adiabatic basis, $H_{\symrm{el}}$ is diagonal, so electronic population transfer between states is mediated entirely by $\sigma_{IJ}$:
\begin{equation}\label{eq:sh_tdse_coefs}
\i\dot c_I=E_I c_I-i\sum_J\sigma_{IJ}c_J.
\end{equation}
With $\symbf{c}(t)=\left(c_1(t),c_2(t),\dots,c_N(t)\right)^{\symrm{T}}$ the Eq~\eqref{eq:sh_tdse_coefs} is a Schr\"odinger equation in a finite-dimensional Hilbert space $\mathbb{C}^N$ and can be rewritten as
\begin{equation}
\i\symbf{\dot{c}}(t)=H_{\symrm{eff}}\symbf{c}(t),\qquad H_{IJ}^{\symrm{eff}}=E_I\delta_{IJ}-\i\sigma_{IJ}
\end{equation}
For electronic dynamics in Eq~\eqref{eq:sh_tdse_coefs} to be unitary, the effective Hamiltonian $H_{\symrm{eff}}$ must be self-adjoint, which in turn requires $\sigma_{IJ}$ to be antisymmetric, $\sigma^\dagger=-\sigma$, in accordance with Sec~\ref{sec:self_adjoint_operators_and_observables}. To evaluate $\sigma_{IJ}$ in practice, note that each adiabatic eigenfunction depends parametrically on the instantaneous nuclear coordinates $\symbf{R}(t)$. By the chain rule,
\begin{equation}
\sigma_{IJ}=\Braket{\phi_I\vert\frac{\partial\phi_J}{\partial t}}=\dot{\symbf{R}}\cdot\Braket{\phi_I\vert\frac{\partial\phi_J}{\partial\symbf{R}}}=\symbf{v}\cdot\symbf{d}_{IJ},
\end{equation}
where $\symbf{v}=\dot{\symbf{R}}$ is the nuclear velocity and $\symbf{d}_{IJ}$ is the \acrfull{nacv}. Although many quantum‐chemistry packages provide \acrshort{nacv} and thus allow computation of the \acrshort{tdc} via $\sigma_{IJ}=\symbf{v}\cdot\symbf{d}_{IJ}$, \acrshort{nacv} can become ill-defined or numerically unstable near conical intersections (e.g., in multi-reference methods). In cases where a reliable \acrshort{nacv} is not available, one may instead employ a direct \acrshort{tdc} approximation using wavefunction overlaps.
